% Доклад по статье «Базис свободных обобщённых алгебр Герстенхабера» (И. Кайгородов, К. Вяткина).
% ≈ 40 минут. Собирается XeLaTeX (или tectonic): tectonic Gerstenhaber_doklad.tex
\RequirePackage{fix-cm}
\documentclass[aspectratio=169,10pt,xcolor=table]{beamer}

\usefonttheme{professionalfonts}
\usepackage{fontspec}
\setsansfont{FiraSans}[Extension=.otf, UprightFont=*-Regular, BoldFont=*-SemiBold,
  ItalicFont=*-Italic, BoldItalicFont=*-SemiBoldItalic]
\renewcommand{\familydefault}{\sfdefault}
\usepackage[russian]{babel}
\usepackage{amsmath,amssymb}
\usepackage{booktabs,array}
\newcolumntype{L}[1]{>{\raggedright\arraybackslash\hyphenpenalty=10000}p{#1}}
\usepackage{tikz}
\usetikzlibrary{positioning,arrows.meta,calc}
\usepackage[most]{tcolorbox}

% ---------------------------------------------------------------- цвета
\definecolor{applegreen}{rgb}{0.55,0.71,0.0}   % #8DB600
\colorlet{deepgreen}{applegreen!58!black}
\colorlet{softgreen}{applegreen!9!white}
\colorlet{midgreen}{applegreen!35!white}
\definecolor{ink}{HTML}{22261C}
\definecolor{mute}{HTML}{70766A}

% ---------------------------------------------------------------- обозначения статьи
\newcommand{\F}{\mathbb F}
\newcommand{\Z}{\mathbb Z}
\newcommand{\G}{\mathfrak G}
\newcommand{\LL}{\mathfrak L}
\newcommand{\D}{\mathfrak D}
\newcommand{\Sm}{\mathcal S}
\newcommand{\Lh}{\widehat{\mathfrak L}}
\newcommand{\rp}[1]{\|#1\|}
\DeclareMathOperator{\spn}{span}
\newcommand{\hl}[1]{{\color{deepgreen}\textbf{#1}}}
\newcommand{\nohyph}{\hyphenpenalty=10000\exhyphenpenalty=10000\relax}

% ---------------------------------------------------------------- оформление
\setbeamertemplate{navigation symbols}{}
\setbeamersize{text margin left=1.25cm, text margin right=1.25cm}
\setbeamercolor{normal text}{fg=ink,bg=white}
\setbeamercolor{structure}{fg=deepgreen}
\setbeamercolor{alerted text}{fg=deepgreen}
\setbeamerfont{frametitle}{size=\Large,series=\bfseries}
\setbeamerfont{footline}{size=\scriptsize}

\setbeamertemplate{background}{%
  \begin{tikzpicture}[remember picture,overlay]
    \fill[applegreen] (current page.north west) rectangle ([xshift=0.22cm]current page.south west);
  \end{tikzpicture}}

\setbeamertemplate{frametitle}{%
  \nointerlineskip\vskip0.4cm
  \begin{beamercolorbox}[wd=\paperwidth,leftskip=1.25cm,rightskip=1.25cm plus 1fil]{frametitle}
    \usebeamerfont{frametitle}\color{ink}\insertframetitle\par
    \vskip0.16cm{\color{applegreen}\rule{2.2cm}{2.2pt}}
  \end{beamercolorbox}\vskip0.02cm}

\setbeamertemplate{footline}{%
  \begin{beamercolorbox}[wd=\paperwidth,ht=2.4ex,dp=1.6ex,leftskip=1.25cm,rightskip=0.9cm]{footline}
    \usebeamerfont{footline}\color{mute}\insertshorttitle\hfill
    {\color{deepgreen}\insertframenumber}\,/\,\inserttotalframenumber
  \end{beamercolorbox}}

\setbeamertemplate{itemize item}{\color{applegreen}\raisebox{0.18ex}{\rule{0.42em}{0.42em}}}
\setbeamertemplate{itemize subitem}{\color{applegreen}\textbullet}
\setbeamertemplate{enumerate item}{\color{deepgreen}\textbf{\insertenumlabel.}}
\setlength{\leftmargini}{1.3em}

% рамки: определение, теорема, пример
\newtcolorbox{defbox}[1]{enhanced, colback=softgreen, colframe=applegreen, boxrule=0pt,
  leftrule=2.8pt, arc=1.5pt, outer arc=1.5pt, left=8pt, right=8pt, top=4pt, bottom=5pt,
  title={#1}, coltitle=deepgreen, colbacktitle=softgreen, fonttitle=\bfseries\small,
  toptitle=4pt, bottomtitle=-2pt, titlerule=0pt}
\newtcolorbox{thmbox}[1]{enhanced, colback=white, colframe=applegreen, boxrule=0.9pt,
  arc=1.5pt, outer arc=1.5pt, left=8pt, right=8pt, top=4pt, bottom=5pt,
  title={#1}, coltitle=white, colbacktitle=applegreen!88!black, fonttitle=\bfseries\small,
  toptitle=2pt, bottomtitle=2pt}
\newtcolorbox{exbox}[1]{enhanced, colback=black!3, colframe=mute!60, boxrule=0pt,
  leftrule=2.8pt, arc=1.5pt, outer arc=1.5pt, left=8pt, right=8pt, top=4pt, bottom=5pt,
  title={#1}, coltitle=mute, colbacktitle=black!3, fonttitle=\bfseries\small,
  toptitle=4pt, bottomtitle=-2pt, titlerule=0pt}

% слайд-разделитель части
\newcommand{\partslide}[2]{%
  {\setbeamertemplate{background}{%
     \begin{tikzpicture}[remember picture,overlay]
       \fill[applegreen] (current page.north west) rectangle (current page.south east);
       \node[anchor=south west, text=white, opacity=0.35, font=\fontsize{96}{96}\selectfont\bfseries, inner sep=0pt]
         at ([xshift=1.25cm,yshift=0.35cm]current page.west) {#1};
       \node[anchor=north west, text=white, font=\fontsize{28}{34}\selectfont\bfseries, text width=0.78\paperwidth, inner sep=0pt]
         at ([xshift=1.3cm,yshift=-0.2cm]current page.west) {\hyphenpenalty=10000 #2};
     \end{tikzpicture}}%
   \setbeamertemplate{footline}{}%
   \begin{frame}[plain,noframenumbering]\end{frame}}}

\title{Базис свободных обобщённых алгебр Герстенхабера}
\author{И.\,Кайгородов, К.\,Вяткина}
\date{}

\begin{document}

% ================================================================ титул
{\setbeamertemplate{background}{}\setbeamertemplate{footline}{}
\begin{frame}[plain]
\begin{tikzpicture}[remember picture,overlay]
  \fill[applegreen] (current page.north west) rectangle ([yshift=-0.56\paperheight]current page.north east);
  \node[anchor=north east, text=white, opacity=0.2, font=\fontsize{40}{40}\selectfont]
    at ([xshift=-0.8cm,yshift=-0.45cm]current page.north east) {$\{1,1\}\neq0$};
  \node[anchor=south west, text width=0.86\paperwidth, align=left, text=white,
        font=\fontsize{23}{29}\selectfont\bfseries]
    at ([xshift=1.25cm,yshift=-0.56\paperheight+0.6cm]current page.north west)
    {\hyphenpenalty=10000\exhyphenpenalty=10000 Базис свободных обобщённых\\ алгебр Герстенхабера};
  \node[anchor=north west, text width=0.82\paperwidth, align=left, text=ink]
    at ([xshift=1.25cm,yshift=-0.56\paperheight-0.75cm]current page.north west)
    {{\Large И.\,Кайгородов, \ К.\,Вяткина}\\[0.55em]
     {\color{mute}Семинар по неассоциативным алгебрам}};
  \fill[applegreen] ([xshift=1.25cm,yshift=-0.56\paperheight-2.75cm]current page.north west)
    rectangle ++(2.2cm,2.2pt);
\end{tikzpicture}
\end{frame}}

% ================================================================ план
\begin{frame}{План доклада}
\vskip0.3cm
\begin{enumerate}\setlength{\itemsep}{0.75em}
  \item \textbf{Контекст:} базисы свободных алгебр, алгебры Пуассона и Герстенхабера
  \item \textbf{Обобщённые алгебры Герстенхабера} и элемент $\{1,1\}$
  \item \textbf{Главная теорема:} базис $U$ свободной алгебры
  \item \textbf{Доказательство:} метод Ширшова, модельная алгебра, якобиатор
  \item \textbf{Размерность} мультилинейной части: $n\cdot n!$
\end{enumerate}
\end{frame}

% ================================================================ 1. контекст
\partslide{01}{Контекст}

\begin{frame}{Базисы свободных алгебр}
\footnotesize\renewcommand{\arraystretch}{1.45}
\begin{tabular}{@{}L{4.2cm}L{5.0cm}L{3.3cm}@{}}
\toprule
\textbf{Класс} & \textbf{Базис свободной алгебры} & \textbf{Кто}\\
\midrule
алгебры Ли & хорошие слова (коммутаторы Холла) & Холл 1950, Ширшов 1958\\
супералгебры Ли & хорошие слова и квадраты нечётных & Штерн 1986; БМПЗ\\
(супер)алгебры Пуассона & одночлены от хороших слов & Шестаков 1993\\
обобщённые супералгебры Пуассона & одночлены от $M\setminus\{1\}$, буква $1$ & Кайгородов 2017\\
\rowcolor{softgreen}
\textbf{обобщённые алгебры Герстенхабера} & \textbf{?} & \textbf{эта работа}\\
\bottomrule
\end{tabular}

\vskip0.45cm
\begin{itemize}\small
  \item Базис свободной алгебры — нормальная форма: тождества, размерности, сравнение классов.
  \item Метод Ширшова (как в K17): \hl{порождение} + \hl{модель}.
\end{itemize}
\end{frame}

\begin{frame}{Соглашения}
\begin{itemize}\setlength{\itemsep}{0.55em}
  \item $\F$ — поле характеристики $0$; все пространства $\Z_2$-градуированы: $V=V_0\oplus V_1$.
  \item $|a|\in\Z_2$ — чётность однородного $a$; \quad \hl{обращённая чётность} $\rp a=|a|+1$.
  \item Суперкоммутативность: $\quad ab=(-1)^{|a||b|}\,ba$.
  \item \hl{Нечётная скобка:} $\quad |\{a,b\}|=|a|+|b|+1 \iff \rp{\{a,b\}}=\rp a+\rp b$.
\end{itemize}
\vskip0.3cm
\begin{defbox}{Нечётная супералгебра Ли}
$$\{a,b\}=-(-1)^{\rp a\rp b}\{b,a\},\qquad
  \{a,\{b,c\}\}=\{\{a,b\},c\}+(-1)^{\rp a\rp b}\{b,\{a,c\}\}.$$
Это ровно супералгебра Ли на $\Pi V$ (та же $V$ с обращённой чётностью).
\end{defbox}
\end{frame}

\begin{frame}{Алгебры Герстенхабера}
\begin{defbox}{Алгебра Герстенхабера (нечётная супералгебра Пуассона)}
Ассоциативная суперкоммутативная супералгебра с $1$ и нечётной скобкой, для которой $\Pi V$ — супералгебра Ли и
$$\{a,bc\}=\{a,b\}c+(-1)^{\rp a|b|}\,b\{a,c\},$$
т.\,е. $\{a,\cdot\}$ — супердифференцирование чётности $\rp a$.
\end{defbox}
\vskip0.15cm
\begin{columns}[T]
\begin{column}{0.5\textwidth}
\textbf{Источники}
\begin{itemize}\small
  \item когомологии Хохшильда $HH^\bullet(A,A)$ (Герстенхабер, 1963)
  \item $\Lambda(\mathfrak g)$ со скобкой Схоутена
  \item BV-алгебры (Баталин–Вилковыский)
\end{itemize}
\end{column}
\begin{column}{0.5\textwidth}
\textbf{Поливекторные поля} $PO(n,n)=\F[x_i]\otimes\Lambda(\xi_i)$:
\small
$$\{f,g\}=\sum_i\Bigl(\tfrac{\partial f}{\partial x_i}\tfrac{\partial g}{\partial \xi_i}+(-1)^{|f|}\tfrac{\partial f}{\partial \xi_i}\tfrac{\partial g}{\partial x_i}\Bigr)$$
\end{column}
\end{columns}
\end{frame}

\begin{frame}{Обобщённые супералгебры Пуассона (чётная скобка)}
\begin{defbox}{Правило Лейбница с поправкой}
$$\{a,bc\}=\{a,b\}c+(-1)^{|a||b|}\,b\{a,c\}-\D(a)\,bc,\qquad \D(a)=\{a,1\}.$$
\end{defbox}
\vskip0.1cm
\begin{exbox}{Пример}
$\F[x]$ со скобкой векторных полей на прямой $\{f,g\}=fg'-f'g$. Здесь $\{f,1\}=-f'$:
$$\{f,gh\}=\{f,g\}h+g\{f,h\}-\{f,1\}\,gh .$$
\end{exbox}
\begin{itemize}\small
  \item Происхождение: двойник Кантора, йордановы скобки (Кантор 1990; Кинг–МакКриммон 1992).
  \item Простые линейно компактные — Кантарини–Кац 2007; базис свободных — Кайгородов 2017.
  \item В чётном случае $\{1,1\}=-\{1,1\}=0$ \ по кососимметричности.
\end{itemize}
\end{frame}

% ================================================================ 2. определения
\partslide{02}{Обобщённые алгебры\\ Герстенхабера}

\begin{frame}{Определение}
\begin{defbox}{Определение 1}
$\Z_2$-градуированное пространство $V$ с $1\in V_0$ и операциями $ab$, $\{a,b\}$:
$$\begin{aligned}
&(1)\quad |ab|=|a|+|b|,\qquad |\{a,b\}|=|a|+|b|+1;\\
&(2)\quad (ab)c=a(bc),\qquad ab=(-1)^{|a||b|}ba,\qquad 1a=a;\\
&(3)\quad \{a,b\}=-(-1)^{\rp a\rp b}\{b,a\};\\
&(4)\quad \{a,\{b,c\}\}=\{\{a,b\},c\}+(-1)^{\rp a\rp b}\{b,\{a,c\}\};\\
&(5)\quad \{a,bc\}=\{a,b\}c+(-1)^{\rp a|b|}b\{a,c\}+(-1)^{\rp a}\D(a)\,bc,\qquad \D(a)=\{1,a\}.
\end{aligned}$$
\end{defbox}
\begin{itemize}\small
  \item $\D=0$ $\Rightarrow$ алгебры Герстенхабера.
  \item Это в точности \emph{odd generalized Poisson superalgebras} Кантарини–Каца (2010).
  \item Естественные примеры с $\D\neq0$: супералгебры $PO(n,n+1)$ (Кантарини–Кац).
\end{itemize}
\end{frame}

\begin{frame}{Что меняет поправка}
\begin{itemize}\setlength{\itemsep}{0.5em}
  \item {\color{mute}Зам. 2.} \ $(1),(3),(4)$ $\iff$ $(\Pi V,\{\,,\})$ — супералгебра Ли.
  \item {\color{mute}Зам. 3.} \ $(5)$ $\iff$ $d_a(b)=\{a,b\}+(-1)^{\rp a}\D(a)\,b$ — супердифференцирование чётности $\rp a$.
  \item Подстановка $a=1$ в $(5)$:
  $$\D(bc)=\D(b)\,c+(-1)^{|b|}\,b\,\D(c)-\D(1)\,bc .$$
\end{itemize}
\begin{thmbox}{Вывод}
$\D$ — нечётное дифференцирование $\iff$ $\D(1)=\{1,1\}=0$.\\[0.3em]
Но $(3)$ при $a=b=1$ даёт $\{1,1\}=\{1,1\}$: \ $\rp1=1$, условия \hl{нет}.
\end{thmbox}
\end{frame}

\begin{frame}{Пример: $\{1,1\}\neq0$}
\begin{exbox}{Пример 1}
$V=\F1\oplus\F\varepsilon$, \ $|\varepsilon|=1$, \ $\varepsilon^2=0$;\qquad
$\{1,1\}=\varepsilon,\quad \{1,\varepsilon\}=\{\varepsilon,1\}=\{\varepsilon,\varepsilon\}=0 .$
\end{exbox}
\begin{itemize}\setlength{\itemsep}{0.5em}
  \item $(3)$ при $a=b=1$: \ $\{1,1\}=\{1,1\}$.
  \item $(4)$: все скобки лежат в $\F\varepsilon$, а $\{\varepsilon,V\}=0$.
  \item $(5)$ при $a=b=c=1$: \quad $\varepsilon=\varepsilon+\varepsilon-\varepsilon$; \ остальные тройки дают $0=0$.
\end{itemize}
\vskip0.2cm
\begin{thmbox}{Итог}
$V$ — обобщённая алгебра Герстенхабера, но $\D$ не дифференцирование:
$$\D(1\cdot1)=\varepsilon\ \neq\ 2\varepsilon=\D(1)\cdot1+1\cdot\D(1).$$
\end{thmbox}
\end{frame}

% ================================================================ 3. базис
\partslide{03}{Свободная алгебра\\ и её базис}

\begin{frame}{Свободная обобщённая алгебра Герстенхабера}
\begin{itemize}\setlength{\itemsep}{0.55em}
  \item $X=\{1<x_1<x_2<\cdots\}$, \ $|1|=0$, \ чётности $|x_i|$ произвольны.
  \item $D_X$ — все неассоциативные слова от $X$ с операциями $\cdot$ и $\{\,,\}$.
  \item $I$ — идеал, порождённый тождествами $(1)$–$(5)$; \quad $\G=D_X/I$.
  \item Универсальность: любое сохраняющее чётность $\{x_i\}\to V$ однозначно продолжается до гомоморфизма $\G\to V$, $1\mapsto1$.
\end{itemize}
\vskip0.3cm
\begin{thmbox}{Ключевое наблюдение}
$1$ — одновременно \hl{единица умножения} и \hl{буква скобки}:
$\{1,x_i\}$, $\{1,1\}$, $\{\{x_i,1\},1\},\ \dots$ — новые элементы.
\end{thmbox}
\end{frame}

\begin{frame}{Хорошие слова (базисные коммутаторы Холла)}
\begin{defbox}{Определение}
Буквы — хорошие слова степени 1. Слово $w$ степени $n$ \hl{хорошее}, если
\begin{enumerate}\small
  \item $w=\{u,v\}$, где $u,v$ — хорошие и $u>v$;
  \item если $u=\{u_1,u_2\}$, то $u_2\le v$.
\end{enumerate}
Порядок: большая степень больше; одна степень — лексикографически по $(u,v)$.
\end{defbox}
\begin{columns}[T]
\begin{column}{0.5\textwidth}\small
\textbf{Примеры} над $\{1<x_1<x_2\}$:
$$\{x_1,1\},\ \{x_2,1\},\ \{x_2,x_1\},$$
$$\{\{x_1,1\},1\},\ \{\{x_1,1\},x_2\},\ \{\{x_2,1\},x_1\}$$
\end{column}
\begin{column}{0.5\textwidth}\small
\textbf{Не хорошее:} $\{\{x_2,x_1\},1\}$, так как $x_1>1$.
По Якоби оно выражается через $\{\{x_2,1\},x_1\}$ и $\{\{x_1,1\},x_2\}$.
\end{column}
\end{columns}
\end{frame}

\begin{frame}{Множество $M$}
$$M=\bigl\{\,u,\ \{v,v\}\ \bigm|\ u,v\ \text{— хорошие слова},\ |v|=0\,\bigr\},\qquad e_1=1<e_2<e_3<\cdots$$
\begin{thmbox}{Лемма 4}
$M$ — базис свободной нечётной супералгебры Ли $\LL$, порождённой $X$.
\end{thmbox}
\begin{itemize}\setlength{\itemsep}{0.45em}
  \item По замечанию 2: $\LL$ — свободная супералгебра Ли на $X$ с чётностью $\rp x$.
  \item Её базис: хорошие слова и квадраты слов с $\rp v=1$ \ (Штерн 1986; БМПЗ; Бокуть–Канг–Ли–Малкольмсон).
  \item $\rp v=1\iff|v|=0$: квадраты \hl{чётных} слов. В частности, $\{1,1\}\in M$.
\end{itemize}
\end{frame}

\begin{frame}{Главная теорема}
$$U=\bigl\{\,e_{i_1}^{k_1}\cdots e_{i_r}^{k_r}\ \bigm|\ r\ge0,\ 1<i_1<\dots<i_r,\ k_j\ge1,\ k_j=1\ \text{при}\ |e_{i_j}|=1\,\bigr\}$$
\centerline{\small\color{mute}суперкоммутативные одночлены от $M\setminus\{1\}$}
\vskip0.25cm
\begin{thmbox}{Теорема 12}
$U$ — базис свободной унитальной обобщённой алгебры Герстенхабера $\G$, порождённой $X$.
Более того, $\G\cong\Sm$ — свободная суперкоммутативная супералгебра над $\LL'=\spn(M\setminus\{1\})$
с явно заданной скобкой.
\end{thmbox}
\begin{itemize}\small
  \item Шаблон тот же, что у Шестакова: «свободная $=$ симметрическая алгебра над свободной скобочной частью».
  \item Новое: буква $1$ в скобках и базисный элемент $\{1,1\}$.
\end{itemize}
\end{frame}

\begin{frame}{Следствие и пример}
\begin{thmbox}{Следствие 13}
В $\G$: \ $\{1,1\}\neq0$, и $\D$ не является дифференцированием.
Более того, $\{1,1\}$, \ $\{1,1\}x_i$, \ $\{1,x_i\}$ \ линейно независимы.
\end{thmbox}
{\small\color{mute}Без теоремы 12: гомоморфизм $\G\to V$ (пример 1), $1\mapsto1$, $x_i\mapsto0$, переводит $\{1,1\}\mapsto\varepsilon\neq0$.}
\vskip0.3cm
\begin{exbox}{Пример 2: мультилинейные элементы $U$ по $1,x_1,\dots,x_n$}
$n=1$: \quad $\{x_1,1\}$\\[0.3em]
$n=2$: \quad $\{x_1,1\}\,x_2,\quad \{x_2,1\}\,x_1,\quad \{\{x_1,1\},x_2\},\quad \{\{x_2,1\},x_1\}$
\end{exbox}
\end{frame}

% ================================================================ 4. доказательство
\partslide{04}{Доказательство}

\begin{frame}{Схема: метод Ширшова}
\centering
\begin{tikzpicture}[
  box/.style={draw=applegreen, line width=0.9pt, fill=softgreen, rounded corners=2pt,
              text width=5.6cm, align=left, inner sep=7pt, font=\small\nohyph},
  arr/.style={-{Stealth[length=7pt]}, line width=1pt, deepgreen}]
  \node[box] (span) {\textbf{\color{deepgreen}Порождение} (лемма 5)\\[0.2em]
     $U$ линейно порождает $\G$: скобки переписываются по $(5)$, индукция по $e$-степени.};
  \node[box, right=1.1cm of span] (model) {\textbf{\color{deepgreen}Модель} (теорема 11)\\[0.2em]
     На $\Sm=S(\LL')$ задаём скобку $\{\,,\}^*$ формулами (7), (8) и проверяем \mbox{аксиомы $(1)$–$(5)$}.};
  \node[box, text width=12.2cm, below=1.0cm of $(span.south)!0.5!(model.south)$, fill=white] (iso)
     {\textbf{\color{deepgreen}Изоморфизм} (теорема 12)\quad
      $\varphi\colon\G\to\Sm$, $\varphi|_X=\mathrm{id}$, переводит $U$ в базис $\Sm$ $\Rightarrow$ $U$ независимо, $\G\cong\Sm$.};
  \draw[arr] (span.south) -- (span.south |- iso.north);
  \draw[arr] (model.south) -- (model.south |- iso.north);
\end{tikzpicture}
\vskip0.45cm
\raggedright\small
Трудная часть — модель: нужна скобка, для которой проверяются все аксиомы, в том числе с $\{1,1\}\neq0$.
\end{frame}

\begin{frame}{Лемма 5: $U$ порождает $\G$}
\begin{itemize}\setlength{\itemsep}{0.45em}
  \item $\psi\colon\LL\to\G$ — гомоморфизм нечётных супералгебр Ли; $\psi(\LL)=\spn M$ (лемма 4).
  \item $\{u,v\}\in\spn U$ для $u,v\in U$ — индукция по $N=\deg_e u+\deg_e v$:
  $$\{u,v_1v_2\}=\{u,v_1\}v_2+(-1)^{\rp u|v_1|}v_1\{u,v_2\}+(-1)^{\rp u}\{1,u\}\,v_1v_2 .$$
  \item $\spn U$ замкнуто относительно обеих операций и содержит $X$ $\Rightarrow$ $\spn U=\G$.
\end{itemize}
\vskip0.2cm
\begin{exbox}{Почему не по числу букв}
$1a=a$ уменьшает число букв, а поправка в $(5)$ увеличивает:
$\{x_1,x_2x_3\}$ (4 буквы) содержит $(-1)^{\rp{x_1}}\{1,x_1\}x_2x_3$ (5 букв).
\end{exbox}
\end{frame}

\begin{frame}{Модельная алгебра}
\small
$\LL'=\spn(M\setminus\{1\})$, \ $\Sm=S(\LL')$, \ $\Lh=\F1\oplus\LL'\cong\LL$ (скобка переносится); \
$\D(u)=\{1,u\}$ для $u\in\Lh$.\\[0.2em]
Для $a=a_1\cdots a_m$: \ $a_{<i}=a_1\cdots a_{i-1}$, \ $a_{>i}=a_{i+1}\cdots a_m$, \ $a_{\hat\imath}=a_{<i}a_{>i}$.
\begin{defbox}{Формулы (7), (8)}
$$\D(a)=\sum_{i=1}^m(-1)^{|a_{<i}|}\,a_{<i}\{1,a_i\}a_{>i}\;-\;(m-1)\{1,1\}\,a,$$
$$\begin{aligned}
\{a,b\}^*=\sum_{i,j}(-1)^{|a_i||a_{>i}|+|b_j||b_{<j}|}\,a_{\hat\imath}\{a_i,b_j\}b_{\hat\jmath}
&-(m-1)\,a\D(b)+(n-1)(-1)^{\rp a}\D(a)\,b\\
&-(m-1)(n-1)\,a\D(1)\,b .
\end{aligned}$$
\end{defbox}
\begin{itemize}
  \item Двойная сумма: «$a_i$ — в правый конец $a$, $b_j$ — в левый конец $b$, скобка соседей».
  \item Остальное — поправки с $\D$; при $\D=0$ остаётся обычное бидифференцирование.
\end{itemize}
\end{frame}

\begin{frame}{Свойства скобки $\{\,,\}^*$}
\begin{thmbox}{Лемма 6}
(7), (8) не зависят от разложения на множители из $\Lh$, и
\begin{enumerate}\small
  \item[(i)] $\{u,v\}^*=\{u,v\}$, \ $\D(u)=\{1,u\}$ при $u,v\in\Lh$;\qquad
  (ii) $\{1,a\}^*=\D(a)$;
  \item[(iii)] $\D(bc)=\D(b)c+(-1)^{|b|}b\D(c)-\D(1)bc$;\qquad
  (iv) $|\{a,b\}^*|=|a|+|b|+1$.
\end{enumerate}
\end{thmbox}
\vskip0.1cm
\begin{columns}[T]
\begin{column}{0.48\textwidth}
\begin{thmbox}{Лемма 7}
$\{a,b\}^*=-(-1)^{\rp a\rp b}\{b,a\}^*$
\end{thmbox}
\end{column}
\begin{column}{0.48\textwidth}
\begin{thmbox}{Лемма 8}
Правило $(5)$ для $\{\,,\}^*$
\end{thmbox}
\end{column}
\end{columns}
\vskip0.1cm
{\small Проверки прямые: формула (8) — замкнутая, без рекурсии.}
\end{frame}

\begin{frame}{Якобиатор}
\begin{thmbox}{Лемма 9}
Пусть выполнены $(1),(2),(3),(5)$ (без тождества Якоби), и
$$J(a,b,c)=\{a,\{b,c\}\}-\{\{a,b\},c\}-(-1)^{\rp a\rp b}\{b,\{a,c\}\}.$$
Тогда
\begin{enumerate}
  \item[(i)] $J(a,b,c)=-(-1)^{\rp a\rp b}J(b,a,c)=-(-1)^{\rp b\rp c}J(a,c,b)$;
  \item[(ii)] $J(a,b,xy)=J(a,b,x)\,y+(-1)^{(|a|+|b|)|x|}\,x\,J(a,b,y)\;{\color{deepgreen}-\,J(a,b,1)\,xy}$.
\end{enumerate}
\end{thmbox}
\begin{itemize}\small
  \item Якобиатор — «дифференцирование с поправкой»; без поправки это классический факт для алгебр Пуассона.
  \item Доказательство: раскрыть три члена по $(5)$, лишние члены сокращаются попарно.
  \item Общая формула, может быть полезна сама по себе.
\end{itemize}
\end{frame}

\begin{frame}{Тождество Якоби и завершение}
\begin{enumerate}\setlength{\itemsep}{0.5em}
  \item \textbf{Предложение 10.} $J=0$ на $\Sm$:
  \begin{itemize}\small
    \item $J(u,v,w)=0$ при $u,v,w\in\Lh$ — это Якоби в $\LL$;
    \item по (ii) и индукции по числу множителей — для любого третьего аргумента;
    \item по (i) — для всех аргументов.
  \end{itemize}
  \item \textbf{Теорема 11.} $(\Sm,\cdot,\{\,,\}^*)$ — унитальная обобщённая алгебра Герстенхабера, порождённая $X$.
  \item \textbf{Теорема 12.} $\varphi\colon\G\to\Sm$ сюръективен, $\varphi(u)=u$ для $u\in U$, а $U$ — базис $\Sm$;
  вместе с леммой 5: $U$ — базис $\G$ и $\varphi$ — изоморфизм. \hfill$\square$
\end{enumerate}
\end{frame}

% ================================================================ 5. размерность
\partslide{05}{Размерность}

\begin{frame}{Мультилинейная часть}
\begin{itemize}\setlength{\itemsep}{0.45em}
  \item $X=\{1,x_1,\dots,x_n\}$; \ $\mathrm{GenGer}(n)$ — оболочка элементов $U$, где каждая из букв $1,x_1,\dots,x_n$ встречается ровно один раз.
  \item Это не однородная компонента: $1a=a$ и поправка в $(5)$ не однородны по $1$.
\end{itemize}
\vskip0.2cm
\begin{thmbox}{Теорема 14}
\centering\Large $\dim\mathrm{GenGer}(n)=n\cdot n!$
\end{thmbox}
\vskip0.2cm
\centering
\begin{tabular}{@{}lcccc@{}}
\toprule
$n$ & 1 & 2 & 3 & 4\\
\midrule
обобщённые алгебры Герстенхабера: $n\cdot n!$ & \textbf{1} & \textbf{4} & \textbf{18} & \textbf{96}\\
алгебры Герстенхабера ($\D=0$, по $x_1,\dots,x_n$): $n!$ & 1 & 2 & 6 & 24\\
\bottomrule
\end{tabular}
\vskip0.15cm
{\small\color{mute}Ответ $n\cdot n!$ совпадает с ответом для обобщённых супералгебр Пуассона (K17).}
\end{frame}

\begin{frame}{Доказательство: перестановки и циклы}
\begin{itemize}\setlength{\itemsep}{0.5em}
  \item Мультилинейный элемент $U$ $\leftrightarrow$ разбиение $\pi$ множества $X$ на блоки, \hl{без блока $\{1\}$}, и мультилинейное хорошее слово на каждом блоке.
  \item Хороших мультилинейных слов на $k$ буквах — $(k-1)!$ (размерность для свободной алгебры Ли).
  \item Квадраты $\{v,v\}$ не мультилинейны.
\end{itemize}
$$\dim\mathrm{GenGer}(n)=\sum_{\pi}\ \prod_{B\in\pi}(|B|-1)!$$
\begin{thmbox}{Подсчёт}
Разбиение $Y$ с циклическим порядком на каждом блоке $=$ перестановка $Y$ (блоки — циклы):
\quad $\sum_{\pi}\prod_B(|B|-1)!=|Y|!$.\\[0.3em]
Вычитаем разбиения с блоком $\{1\}$: \quad $(n+1)!-n!=n\cdot n!$.
\end{thmbox}
\end{frame}

% ================================================================ итоги
\begin{frame}{Итоги}
\begin{itemize}\setlength{\itemsep}{0.6em}
  \item Построен базис $U$ свободной обобщённой алгебры Герстенхабера: одночлены от хороших слов и квадратов чётных хороших слов, где $1$ — буква скобки.
  \item Новое явление: $\{1,1\}\neq0$, и $\D$ — не дифференцирование.
  \item Модель задана замкнутой формулой для скобки; Якоби — через общую лемму о якобиаторе.
  \item $\dim\mathrm{GenGer}(n)=n\cdot n!$.
  \item При $\D=0$ получается известный базис свободной алгебры Герстенхабера (замечание 15).
\end{itemize}
\end{frame}

\begin{frame}{Литература}
\footnotesize
\begin{itemize}\setlength{\itemsep}{0.35em}
  \item M.~Hall, \emph{A basis for free Lie rings and higher commutators in free groups}, Proc. AMS \textbf{1} (1950), 575--581.
  \item A.~Shtern, \emph{Free Lie superalgebras}, Sib. Math. J. \textbf{27} (1986), 136--140.
  \item I.~Shestakov, \emph{Quantization of Poisson superalgebras and speciality of Jordan Poisson superalgebras}, Algebra and Logic \textbf{32} (1993), 309--317.
  \item N.~Cantarini, V.~Kac, \emph{Classification of linearly compact simple Jordan and generalized Poisson superalgebras}, J. Algebra \textbf{313} (2007), 100--124.
  \item N.~Cantarini, V.~Kac, \emph{Classification of linearly compact simple rigid superalgebras}, IMRN (2010), 3341--3393.
  \item I.~Kaygorodov, \emph{Algebras of Jordan brackets and generalized Poisson algebras}, Linear Multilinear Algebra \textbf{65} (2017), 1142--1157.
  \item M.~Gerstenhaber, \emph{The cohomology structure of an associative ring}, Ann. Math. \textbf{78} (1963), 267--288.
\end{itemize}
\end{frame}

% ================================================================ финал
{\setbeamertemplate{background}{%
   \begin{tikzpicture}[remember picture,overlay]
     \fill[applegreen] (current page.north west) rectangle (current page.south east);
     \node[text=white, opacity=0.25, font=\fontsize{54}{54}\selectfont, anchor=south east]
       at ([xshift=-0.9cm,yshift=0.8cm]current page.south east) {$\dim=n\cdot n!$};
     \node[text=white, font=\fontsize{34}{40}\selectfont\bfseries, anchor=west]
       at ([xshift=1.3cm,yshift=0.5cm]current page.west) {Спасибо за внимание!};
     \node[text=white, anchor=north west, font=\large]
       at ([xshift=1.3cm,yshift=-0.35cm]current page.west) {Вопросы?};
   \end{tikzpicture}}%
 \setbeamertemplate{footline}{}%
 \begin{frame}[plain,noframenumbering]\end{frame}}

\end{document}
